% !TeX root = ../proyecto.tex

\section{Requerimientos Funcionales}

A continuación se enlistan los requerimientos funcionales del sistema, los cuales definen las
capacidades de recolección, análisis, investigación e interacción con el dashboard web.

\begin{requerimientos}
    % \FRitem{ID}{Nombre}{Descripción}{Prioridad}{Referencia}

    \FRitem{RF-01}{Recolección Multifuente}{El sistema debe extraer noticias a partir de canales
    RSS, consultas directas a Google News y scraping HTML con emulación de fingerprint TLS
    (\texttt{StealthSession} v7.0) sobre fuentes sin feed RSS.}{Muy Alta}{CU-01, CU-02}

    \FRitem{RF-02}{Deduplicación en Cascada}{El sistema debe filtrar noticias repetidas mediante
    un pipeline de cuatro fases secuenciales: Hash MD5 exacto, similitud de Jaccard sobre tokens
    ($J \geq 0.70$), SimHash (distancia Hamming $\leq 6$ bits) y similitud coseno TF-IDF
    ($\cos \geq 0.85$).}{Alta}{CU-01}

    \FRitem{RF-03}{Clasificación Escalonada (Tiered Scoring v7.1)}{El motor semántico debe procesar
    las noticias a través de 4 capas autónomas: Capa~0 Escudo Léxico (bloqueo absoluto por dominio
    ajeno), Capa~1 Escudo Suave (penalización aditiva $\delta = -0.35$ por contenido legislativo),
    Capa~2 Inferencia BETO Fine-Tuned (o Zero-Shot como fallback) y Capa~3 Bonificación Heurística
    por evidencia explícita de orfandad.}{Muy Alta}{CU-01}

    \FRitem{RF-04}{Agrupamiento Semántico (BERTopic)}{El sistema debe emplear el pipeline
    UMAP + HDBSCAN + c-TF-IDF de BERTopic para agrupar noticias sobre el mismo hecho fáctico
    en tópicos semánticos, con reducción de outliers de 81.4\% a máximo 55\% mediante tres
    estrategias progresivas.}{Alta}{CU-01}

    \FRitem{RF-05}{Investigación Profunda Bajo Demanda}{El sistema debe sintetizar una ficha de
    caso estructurada en JSON con siete campos (\texttt{ubicacion}, \texttt{victimas},
    \texttt{ninos\_afectados}, \texttt{edades}, \texttt{situacion\_actual},
    \texttt{medios\_contacto}, \texttt{resumen}) a partir de una noticia, usando DuckDuckGo
    (HTML + Lite) y el LLM Gemini Flash, ejecutado de forma asíncrona en hilo de fondo.}{Muy Alta}{CU-03}

    \FRitem{RF-06}{Ingestión Manual por Enlace Externo}{El analista debe poder introducir una URL
    periodística arbitraria para que el pipeline de \texttt{DeepInvestigator} extraiga su contenido
    mediante \texttt{StealthSession}, valide el cuerpo periodístico con Gemini Flash y genere la
    ficha de caso correspondiente.}{Media}{CU-04}

    \FRitem{RF-07}{Gestión de Casos en Seguimiento}{El sistema debe permitir al analista marcar
    noticias investigadas (\texttt{en\_seguimiento = True} en el JSON de investigación) para
    mantenerlas en una lista de monitoreo activo consultable en la pantalla \IUref{IU-04}{Seguimiento}.
    }{Alta}{CU-05}

    \FRitem{RF-08}{Exportación CSV de Seguimiento}{El sistema debe permitir descargar los registros
    en seguimiento en formato tabular (\texttt{seguimiento\_casos\_nna.csv}), aplanando los campos
    del JSON de investigación en columnas individuales (ubicación, víctimas, NNA afectados, edades,
    situación actual, scores semánticos y enlace original).}{Alta}{CU-05}

    \FRitem{RF-09}{Búsqueda de Texto Completo (FTS)}{El sistema debe permitir búsquedas semánticas
    sobre el corpus mediante FTS nativo de PostgreSQL con índices GIN y stemming en español
    (diccionario Snowball), con fallback a búsqueda textual expandida con sinónimos del dominio
    (\texttt{SynonymDictionary}) sobre CSV.}{Alta}{CU-07}

    \FRitem{RF-10}{Exportación CSV del Análisis Completo}{El sistema debe permitir descargar la
    totalidad del corpus clasificado en formato CSV, exportando directamente desde PostgreSQL
    o desde el archivo de datos en memoria como alternativa.}{Media}{CU-01}

    \FRitem{RF-11}{Gráficas y Estadísticas del Corpus}{El sistema debe proveer datos para cuatro
    visualizaciones: distribución temporal de noticias por mes, distribución de relevancia
    (Alta/Media/Baja/No relevante), distribución geográfica por estado de México (top 15) y
    top 10 de fuentes por volumen de noticias.}{Media}{CU-01}

    \FRitem{RF-12}{Gestión del Catálogo de Fuentes}{El administrador debe poder crear, leer,
    actualizar y eliminar fuentes periodísticas en el catálogo. Fuentes predeterminadas
    (\texttt{is\_predefined = True}) solo pueden desactivarse, no eliminarse. El sistema debe
    ofrecer sondeo técnico previo (\texttt{probe\_url}) que detecte tipo de fuente, accesibilidad
    WAF y feeds RSS disponibles sin persistir la URL.}{Alta}{CU-02}

    \FRitem{RF-13}{Autenticación de Usuarios}{El sistema debe permitir el inicio y cierre de
    sesión mediante credenciales (usuario y contraseña), establecer sesiones persistentes con
    Flask-Login y proteger todos los endpoints del dashboard con el decorador
    \texttt{@login\_required}. Las rutas de salud del sistema (\texttt{/health}) quedan exentas.
    }{Alta}{CU-06}

    \FRitem{RF-14}{Limpieza de Historial}{El sistema debe permitir al administrador borrar
    completamente el historial de análisis, incluyendo todos los registros de la base de datos
    PostgreSQL, los archivos CSV del directorio \texttt{data/} y la caché de URLs procesadas
    (\texttt{seen\_urls.json}), para reiniciar el pipeline desde cero.}{Baja}{CU-01}

\end{requerimientos}
\FootnotePrioridad
